% ====================================================================
% graphite_ica_ch2_v4-05.tex — Chapter 2 (v4-05: 코인 하프셀 엔트로피·가역
%   발열의 통계열역학, Opus build)
%   ★v4 = intercalation 엔트로피·가역 발열의 *통계열역학 챕터* — 분포를 명시
%        전개(분배함수 Z → 점유 분포 ⟨n⟩ = Ch1 logistic 의 기원 → configurational
%        ·vibrational(BE)·electronic(FD) 엔트로피 → 섞임/겹침 닫힌식(A·B·C·D)
%        → 극한·코너(I) → 가역 발열 = 분포 재배열 열).
%   범위: 코인 하프셀 단독(단일 전극 vs Li). 전셀 합성(∂U_cell) 범위 외.
%   build: xelatex 2-pass (kotex).
% ====================================================================
\documentclass[11pt,a4paper]{article}
\usepackage[margin=22mm]{geometry}
\usepackage{kotex}
\setmonofont{D2Coding}[Scale=MatchLowercase]
\usepackage{amsmath,amssymb,mathtools,bm}
\usepackage{booktabs,longtable,array}
\usepackage{enumitem}
\usepackage{amsthm}
\usepackage{tikz}
\usetikzlibrary{calc,arrows.meta,decorations.pathreplacing,positioning}
\usepackage{hyperref}
\usepackage{fancyhdr}
\usepackage{lastpage}
\usepackage{setspace}
\setstretch{1.12}
\setlength{\parskip}{0.45em}\setlength{\parindent}{0pt}\setlength{\headheight}{14pt}
\setlength{\emergencystretch}{2em}
\hypersetup{colorlinks=true, linkcolor=blue, urlcolor=blue, citecolor=blue,
  pdftitle={흑연 음극 dQ/dV — Chapter 2 (v4): 코인 하프셀 엔트로피·가역 발열의 통계열역학}}
\pagestyle{fancy}\fancyhf{}
\lhead{흑연 음극 $dQ/dV$ — Ch.2 (v4, 엔트로피 통계열역학)}\rhead{\thepage/\pageref{LastPage}}
\renewcommand{\headrulewidth}{0.4pt}
\newtheorem*{keybox}{핵심}
\newtheorem*{srcbox}{문헌 근거}
\newtheorem*{warnbox}{주의}
\newtheorem*{derbox}{유도}
\newtheorem*{procedurebox}{계산 절차}
\newcommand{\dd}{\mathrm{d}}
\newcommand{\rxn}{\mathrm{rxn}}
\newcommand{\eq}{\mathrm{eq}}
\newcommand{\rev}{\mathrm{rev}}
\newcommand{\irr}{\mathrm{irr}}
\newcommand{\oc}{\mathrm{oc}}
\newcommand{\config}{\mathrm{config}}
\newcommand{\vib}{\mathrm{vib}}
\newcommand{\eff}{\mathrm{eff}}
\newcommand{\hys}{\mathrm{hys}}
\newcommand{\kB}{k_{\mathrm B}}
\newcommand{\avg}[1]{\langle #1\rangle}
\renewcommand{\theequation}{2.\arabic{equation}}
\renewcommand{\thesection}{2.\arabic{section}}
\title{\textbf{\LARGE Chapter 2}\\[0.35em]\Large 코인 하프셀 엔트로피와 가역 발열의 통계열역학\\[0.2em]\large — 분배함수에서 점유 분포로, 분포에서 $\partial U/\partial T\cdot\Delta S(x)\cdot\dot Q_\rev$ 로}
\author{}\date{}

\begin{document}
\maketitle

% ====================================================================
\section*{서 — 왜 엔트로피에는 통계가 필요한가}
\addcontentsline{toc}{section}{서}
% ====================================================================
Chapter 1 은 흑연 음극의 한 $\dd Q/\dd V$ 곡선을, 전이별 평형 중심 전위
$U_j(T)=(-\Delta H_{\rxn,j}+T\,\Delta S_{\rxn,j})/F$ 위에 세운 평형 종과 그 위에 얹힌 동역학 꼬리로
설명했다. 그 곡선의 평형 골격에는 이미 한 가지 ``열''이 잠겨 있다. 평형 중심의 \emph{온도 의존}
$\partial U_j/\partial T$ 인데, 열역학 1$\cdot$2법칙은 이 기울기를 곧장 반응 엔트로피로 환산한다. 그리고
반응 엔트로피야말로 전지가 충방전하며 흡수$\cdot$방출하는 \emph{가역(엔트로피) 발열}의 원천이다. 본 장의
일은 이 잠긴 열을 끄집어내어 정량화하는 것이다.

문제는, 엔트로피를 ``$\partial U/\partial T$ 를 측정하면 나오는 수''로만 다루면 그 \emph{모양}을 이해할 수
없다는 데 있다. 흑연 하프셀의 측정 엔트로피 계수 $\partial U_\oc/\partial T(x)$ 는 SOC 에 따라 부호를 바꾸고
($x$ 작을 때 양, 클 때 음), 전이 경계마다 봉우리와 골을 만들며, 어떤 자리에서는 발산에 가깝게 치솟는다. 이
비선형 모양을 ``전이마다 상수 $\Delta S_{\rxn,j}$ 를 가정''하는 평형 정전 식만으로는 재현하지 못한다. 모양은
\emph{분포}에서 온다 — Li 이온이 격자 자리들에 \emph{어떻게 흩어져 앉아 있는가}, 격자 진동이 포논 모드에
\emph{어떻게 분배되는가}, 전도 전자가 Fermi 준위 부근에 \emph{어떻게 채워지는가}. 엔트로피는 본질적으로
이 미시상태 분포의 흩어짐 척도이고, 가역 발열
\[
\dot Q_\rev=-I\,T\,\frac{\partial U_\oc}{\partial T}=-\frac{I\,T}{F}\,\Delta S(x)
\]
은 한 충전상태에서 다른 충전상태로 넘어갈 때 \emph{이 분포를 재배열하면서 주고받는 열}이다. 곧
$\partial U_\oc/\partial T(x)$ 의 모양을 이해하려면, 그 밑에 깔린 점유 분포를 명시적으로 세워야 한다.

그래서 본 장은 통계열역학 챕터다. 분포를 결론으로 인용하는 대신 \emph{본체}로 전개한다. 출발점은 단일 격자
자리의 분배함수 $Z$ 이고(\S\ref{sec:partition}), 거기서 점유 분포 $\avg{n}$ 을 얻는데 이것이 바로 Chapter 1
의 logistic 평형 점유의 통계역학적 기원임을 보인다. 점유 분포에서 Boltzmann$\cdot$Shannon 엔트로피
$S_\config=-R\sum p\ln p$ 가 나오고(\S\ref{sec:config}), 같은 분포 언어로 격자 진동(포논 Bose--Einstein)과
전도 전자(Fermi--Dirac)의 엔트로피를 더한다(\S\ref{sec:vibel}). 세 분포가 합쳐져 측정 $\partial U_\oc/\partial T(x)$
의 모든 모양 — 봉우리 겹침, 봉우리 내부 발산, 상호작용에 의한 폭 변형, 히스테리시스 분기 — 을 만든다
(\S\ref{sec:mixing}). 극한과 코너에서 검산하고(\S\ref{sec:limits}), 마지막으로 분포 재배열로서의 가역
발열로 닫는다(\S\ref{sec:revheat}). 사슬은
\[
\boxed{\;Z \;\to\; \avg{n}=\text{점유 분포}\;\to\; S=-R\textstyle\sum p\ln p \;\to\; \frac{\partial U_\oc}{\partial T}(x)=\frac{\Delta S(x)}{F}\;\to\;\dot Q_\rev=-I\,T\,\frac{\partial U_\oc}{\partial T}\;}
\]
이고, 각 화살표는 점프 없이 분포에서 유도된다.

\begin{warnbox}
\textbf{범위 — 코인 하프셀 단독.} 본 장은 단일 전극(흑연 vs.\ Li 금속)의 하프셀 $\partial U_\oc/\partial T$ 와
그 가역 발열만 다룬다. 전셀 합성 $\partial U_\mathrm{cell}/\partial T=\partial U_\mathrm{cat}/\partial T-\partial
U_\mathrm{an}/\partial T$ (캐소드$-$애노드 차감)는 \emph{범위 밖}이다 — 본 장의 분포는 한 전극의 점유 분포다.
주 사례는 흑연 음극 하프셀이며, LCO 의 전자 엔트로피 항은 분포 언어의 한 사례로만 포함한다.
\end{warnbox}

% ====================================================================
\section{격자기체 분배함수에서 점유 분포로}\label{sec:partition}
% ====================================================================
통계역학에서 한 계의 모든 평형 열역학량은 분배함수에서 나온다. 흑연 격자에 Li 이온이 삽입되는 문제는
\emph{격자기체}(lattice gas)로 모형화된다 \cite{newman,huggins2009}: 결정 안에 동등한 삽입 자리들이 있고,
각 자리는 비어 있거나(에너지 0) Li 하나가 점유한다(에너지 $\varepsilon_0$). 자리들은 전해질을 통해 Li 저장조
(상대극 Li 금속)와 평형을 이루므로 전기화학 퍼텐셜 $\mu$ 가 고정된 \emph{대정준}(grand-canonical) 문제다.
이 절은 그 분배함수를 세우고, 거기서 점유 분포 $\avg{n}$ 을 끌어낸 뒤, 그것이 Chapter 1 이 \emph{이미 쓰고
있던} logistic 평형 점유의 기원임을 보인다.

\subsection{단일 자리 대정준 분배함수}
자리 하나만 떼어 보자. 두 미시상태 — 빈 자리($n=0$, 에너지 0)와 점유($n=1$, 에너지 $\varepsilon_0$) — 의
대정준 분배함수는 각 상태에 Gibbs 인자 $e^{-\beta(\varepsilon_0-\mu)n}$ 을 곱해 더한 것이다($\beta=1/\kB T$):
\begin{equation}
Z_1 \;=\; \sum_{n=0,1} e^{-\beta(\varepsilon_0-\mu)\,n}
\;=\; 1 \;+\; e^{-\beta(\varepsilon_0-\mu)}.
\label{eq:Z1}
\end{equation}
자리의 평균 점유는 정의상 분배함수의 로그 미분, 곧 두 상태의 확률 가중 평균이다:
\begin{equation}
\avg{n} \;=\; \frac{0\cdot 1 + 1\cdot e^{-\beta(\varepsilon_0-\mu)}}{Z_1}
\;=\; \frac{e^{-\beta(\varepsilon_0-\mu)}}{1+e^{-\beta(\varepsilon_0-\mu)}}
\;=\; \frac{1}{1+e^{\beta(\varepsilon_0-\mu)}}.
\label{eq:occ}
\end{equation}
이 $\avg{n}$ 은 자리가 채워질 확률 $p_1$ 그 자체이고, 빌 확률은 $p_0=1-\avg{n}$ 이다. 식~\eqref{eq:occ} 는
구조상 \emph{Fermi--Dirac 함수}와 동형이다 — 자리당 2-상태(채움/빔)라는 점이 전자 준위의 점유/공위와 같은
대수 구조를 낳는다. 다른 점은 ``입자''가 페르미온 전자가 아니라 Li 이온이고, 단일 준위 에너지가
$\varepsilon_0-\mu$ 라는 것뿐이다.

\subsection{전기화학 퍼텐셜과 Chapter 1 logistic 의 기원}\label{ssec:logistic}
이제 화학을 넣는다. Li 가 자리에 들어가는 전기화학 반응의 전기화학 퍼텐셜은 셀 전위에 선형으로 묶인다:
\begin{equation}
\mu \;=\; \mu^0 \;-\; \sigma_d\,F\,(V-U_j),
\label{eq:muV}
\end{equation}
여기서 $V$ 는 (평형) 전극 전위, $U_j$ 는 그 전이의 표준(중심) 전위, $F$ 는 Faraday 상수, $\sigma_d=\pm1$ 은
충/방전 방향 부호(Chapter 1 규약)다. 전위가 오르면 ($V>U_j$) Li 의 전기화학 퍼텐셜이 낮아져 자리에서
빠져나가므로 \emph{점유} $\avg{n}=\theta$ 는 $V$ 에 따라 \emph{감소}한다. 식~\eqref{eq:muV}($\sigma_d=+1$ 평형
규약)를 \eqref{eq:occ} 에 대입하고 $\varepsilon_0-\mu^0\equiv0$ 을 전이 중심 기준으로 잡으면, 평형 점유와 그
여집합(탈리튬화 진행률 $\xi\equiv1-\theta$)은
\begin{equation}
\boxed{\;\theta_\eq(V,T)=\avg{n}=\frac{1}{1+\exp\!\bigl[+(V-U_j)/w\bigr]},\qquad
\xi_\eq(V,T)=1-\theta_\eq=\frac{1}{1+\exp\!\bigl[-(V-U_j)/w\bigr]}\;,}
\label{eq:logistic}
\end{equation}
이고 여기서 $w\equiv RT/F$ 다(몰 단위 환산 $R=N_A\kB$, $F=N_A e$ 이므로 $\beta(\varepsilon_0-\mu)$ 의 분자를
몰 에너지로 쓰면 지수가 $(V-U_j)/(RT/F)=(V-U_j)/w$ 가 되어 $\beta F=1/w$). 곧 \emph{점유} $\theta$ 는 $V$ 에
감소하는 logistic, \emph{탈리튬화 진행률} $\xi=1-\theta$ 는 $V$ 에 증가하는 logistic이며, 후자가 정확히
Chapter 1 의 전이별 \emph{logistic 평형 종}이다($\xi\to1$ 고전위$\cdot$Li 희박, $\xi\to0$ 저전위$\cdot$만충).

\begin{keybox}
Chapter 1 의 평형 종 $\xi_\eq$ 는 \emph{현상학적 곡선맞춤}이 아니라 \emph{격자기체 점유 분포의 여집합}(탈리튬화
진행률 $=1-\theta$)이다. logistic 폭 $w=RT/F$ 는 임의 모수가 아니라 단일 자리 2-상태 분배함수~\eqref{eq:Z1} 가
정하는 분포의 열적 폭이다. ★표기: $\theta=$ Li 점유율(만충서 1), $\xi=1-\theta=$ 탈리튬화 진행률(희박서 1);
둘은 같은 점유 분포의 두 이름이다. Chapter 2 는 이 분포를 결론이 아니라 출발점으로 삼아 엔트로피를 쌓는다.
\end{keybox}

$\xi=1-\theta$ 와 전위의 관계를 뒤집으면(점유 $\theta$ 의 Nernst 관계 $V=U_j+(RT/F)\ln[(1-\theta)/\theta]$ 에
$\theta=1-\xi$ 대입)
\begin{equation}
V(\xi) \;=\; U_j \;+\; \frac{RT}{F}\,\ln\!\frac{\xi}{1-\xi},
\label{eq:Vxi}
\end{equation}
이 되는데, 이 \emph{Nernst 형 로그항}이 뒤에서 configurational 엔트로피의 부분몰 형태로 다시 나타난다
(\S\ref{sec:config}). 식~\eqref{eq:Vxi} 의 $\ln[\xi/(1-\xi)]$ 가 단순한 곡선맞춤이 아니라 점유 분포의
엔트로피에서 온다는 것이 이 장의 첫 매듭이다.

\subsection{다자리 평균장 — Bragg--Williams 와 상호작용}\label{ssec:BW}
실제 격자에서는 자리들이 독립이 아니다. 점유된 이웃 자리는 다음 Li 의 삽입 에너지를 바꾼다(정전 반발,
탄성 변형, 전자 구조 변화). 이 상호작용을 평균장으로 다루는 것이 Bragg--Williams 근사다 \cite{huggins2009,
bazant2013}. 점유율 $\theta\in[0,1]$ (자리당 평균 점유, 거시적으로 $\avg{n}$ 과 같음) 의 자유에너지는
\begin{equation}
g(\theta) \;=\; g_0 \;+\; RT\bigl[\theta\ln\theta + (1-\theta)\ln(1-\theta)\bigr]
\;+\; \Omega\,\theta(1-\theta),
\label{eq:BW}
\end{equation}
이고, 우변 둘째 항이 \emph{혼합 엔트로피}(곧 \S\ref{sec:config}), 셋째 항이 \emph{혼합 엔탈피}(상호작용
모수 $\Omega$, regular-solution 형)다. $\Omega>0$ 은 같은 종끼리 뭉치려는 인력(상분리 경향), $\Omega<0$ 은
교대 정렬(규칙화) 경향을 뜻한다. 평형 전위는 $\theta$ 1몰 삽입당 자유에너지 변화이므로
\begin{equation}
V_\eq(\theta) \;=\; U_j \;+\; \frac{RT}{F}\ln\!\frac{\theta}{1-\theta} \;+\; \frac{\Omega}{F}\,(1-2\theta),
\label{eq:Veq_BW}
\end{equation}
가 되어, 식~\eqref{eq:Vxi} 의 이상 격자기체에 상호작용 항 $\Omega(1-2\theta)/F$ 가 더해진다. 이 상호작용은
뒤에서 두 가지 결과를 낳는다: (i) 점유 분포(따라서 $\dd Q/\dd V$ 봉우리)의 \emph{폭}을 바꾸고
(\S\ref{ssec:weff}), (ii) $\Omega$ 가 임계값을 넘으면 분포가 쌍봉(bimodal)이 되어 \emph{상분리}가 일어난다.
임계 조건은 $\partial V_\eq/\partial\theta=0$ 의 첫 등장, 곧
\begin{equation}
\frac{\dd V_\eq}{\dd\theta} \;=\; \frac{RT}{F}\frac{1}{\theta(1-\theta)} \;-\; \frac{2\Omega}{F},
\qquad
\theta=\tfrac12\text{ 에서 }\;\frac{\dd V_\eq}{\dd\theta}\Big|_{1/2}=\frac{4RT-2\Omega}{F},
\label{eq:slope_BW}
\end{equation}
이 0 이 되는 $\Omega=2RT$ 다. $\Omega<2RT$ 면 등온선이 단조(균질 고용체), $\Omega>2RT$ 면 중앙에서
기울기가 음으로 뒤집혀(불안정) 상분리한다. 이 임계 $\Omega=2RT$ 가 \S\ref{ssec:weff}$\cdot$\S\ref{sec:limits}
에서 핵심 코너로 다시 나온다.

% ====================================================================
\section{점유 분포에서 configurational 엔트로피로}\label{sec:config}
% ====================================================================
이제 분포에서 엔트로피를 끌어낸다. 한 자리가 점유될 확률 $\theta$, 빌 확률 $1-\theta$ 인 베르누이 분포의
정보 엔트로피(Boltzmann--Shannon 엔트로피)는, 자리 1몰당
\begin{equation}
\boxed{\;S_\config \;=\; -R\sum_i p_i\ln p_i
\;=\; -R\bigl[\theta\ln\theta + (1-\theta)\ln(1-\theta)\bigr]\;.}
\label{eq:Sconfig}
\end{equation}
이것은 Bragg--Williams 자유에너지~\eqref{eq:BW} 의 둘째 항을 $-T$ 로 나눈 것과 정확히 일치한다($g$ 의 엔트로피
부분 $=-T S_\config$). 곧 \S\ref{ssec:BW} 에서 자유에너지로 도입한 혼합 항이, 점유 분포의 엔트로피로서
독립적으로 유도된다 — 자유에너지 모형과 분포 통계가 같은 식으로 만나는 지점이다.

\subsection{부분몰 configurational 엔트로피 — 발산하는 봉우리 내부 항}
가역 발열에 들어가는 것은 \emph{Li 1몰 삽입당} 엔트로피 변화, 곧 부분몰량이다. 식~\eqref{eq:Sconfig} 를
$\theta$ 로 미분하면
\begin{equation}
\frac{\partial S_\config}{\partial\theta}
\;=\; -R\bigl[\ln\theta + 1 - \ln(1-\theta) - 1\bigr]
\;=\; -R\,\ln\!\frac{\theta}{1-\theta},
\label{eq:dSconfig}
\end{equation}
가 된다. 이 한 줄이 이 장의 중심 결과 중 하나다. 식~\eqref{eq:dSconfig} 와 식~\eqref{eq:Vxi} 의 logistic
역함수를 나란히 두면,
\begin{equation}
V(\xi)=U_j+\frac{RT}{F}\ln\!\frac{\xi}{1-\xi}
\;\Longrightarrow\;
\frac{\partial V}{\partial T}\Big|_\xi
= \frac{\Delta S_{\rxn,j}}{F}\;+\;\frac{R}{F}\ln\!\frac{\xi}{1-\xi}
= \frac{\Delta S_{\rxn,j}}{F}\;+\;\frac{1}{F}\frac{\partial S_\config}{\partial\theta}\Big|_{\theta=1-\xi},
\label{eq:dVdT_config}
\end{equation}
임을 확인한다($\partial U_j/\partial T=\Delta S_{\rxn,j}/F$ 를 첫 항으로 쓰고, $w=RT/F$ 의 명시적 $T$ 의존이
둘째 항을 낳음; 점유 $\theta=1-\xi$ 이므로 $\partial S_\config/\partial\theta|_{\theta=1-\xi}=-R\ln[(1-\xi)/\xi]
=+R\ln[\xi/(1-\xi)]$ — Li 삽입은 점유 $\theta$ 를 늘리는 방향이라 부호가 $+$ 로 맞물린다). 곧 \textbf{Chapter 1 의
logistic 폭 $w=RT/F$ 가 이미 부분몰 configurational 엔트로피를 담고 있었다}. 봉우리 내부 $\partial U_\oc/\partial T(x)$
의 $x$-의존(비선형 모양)은 새 물리가 아니라, 점유 분포의 엔트로피가 $w$ 의 온도 의존을 통해 자동으로 들어온 것이다.

부호는 명확하다(그림~\ref{fig:occ_config}; 규약: $\xi$ = 탈리튬화 진행률, $\xi\to1$ = Li 희박, $\xi\to0$ = 만충):
\begin{itemize}[leftmargin=1.4em,itemsep=1pt]
\item $\xi\to1$ (희박, Li-poor): $\ln[\xi/(1-\xi)]\to+\infty$ — config 부분몰 엔트로피 $\to+\infty$. 희박
한계에서 한 Li 를 더 넣을 자리의 \emph{선택지}가 폭증하므로 엔트로피가 발산한다. 이것이 저-$x$ 에서
관측되는 큰 양의 $\Delta S$($+29$ J\,mol$^{-1}$K$^{-1}$ 급)의 주원천이다 — dilute 격자기체 모형이 같은
발산을 예측한다 \cite{occupation2019}.
\item $\xi\to0$ (만충, Li-rich): $\ln[\xi/(1-\xi)]\to-\infty$ — config 부분몰 엔트로피 $\to-\infty$. 거의 다 찬
격자에 마지막 Li 를 끼우면 배열 선택지가 사라져 엔트로피가 음으로 발산한다.
\item $\xi=\tfrac12$ (중심): $\ln 1=0$ — config 기여 $=0$. 봉우리 중심에서 부분몰 config 엔트로피가 정확히
사라지고, $\partial V/\partial T$ 는 표준값 $\Delta S_{\rxn,j}/F$ 로 환원된다.
\end{itemize}

\begin{figure}[t]
\centering
\begin{tikzpicture}[x=1.05cm,y=1.0cm]
% ---- left panel: occupation distribution <n>(eta) ----
\begin{scope}
  \draw[->,>=Stealth] (-2.6,0) -- (2.8,0) node[right]{\small $(\mu-\varepsilon_0)/RT$};
  \draw[->,>=Stealth] (0,-0.15) -- (0,3.0) node[above]{\small $\langle n\rangle$};
  \draw[dashed] (-2.6,2.5) -- (2.6,2.5) node[right,yshift=-2pt]{\scriptsize $1$};
  \draw[dashed] (0,1.25) -- (0.12,1.25); \node[left] at (-0.05,1.25){\scriptsize $\tfrac12$};
  % logistic curve n = 1/(1+e^{-eta}), scaled to height 2.5
  \draw[very thick,blue] plot[smooth,domain=-2.6:2.6,samples=60]
    (\x, {2.5/(1+exp(-2*\x))});
  \node[blue] at (1.7,2.15) {\scriptsize $\langle n\rangle=\dfrac{1}{1+e^{-(\mu-\varepsilon_0)/RT}}$};
  \node[below] at (0,-0.55) {\small (a) 점유 분포 = Ch1 logistic};
\end{scope}
% ---- right panel: S_config(theta) ----
\begin{scope}[xshift=7.3cm]
  \draw[->,>=Stealth] (-0.15,0) -- (3.0,0) node[right]{\small $\theta$};
  \draw[->,>=Stealth] (0,-0.15) -- (0,3.0) node[above]{\small $S_\mathrm{config}/R$};
  \draw[dashed] (2.5,0) -- (2.5,0.1) node[below,yshift=-3pt]{\scriptsize $1$};
  % S = -[theta ln theta + (1-theta) ln(1-theta)], max ln2~0.693 at 0.5; x:theta*2.5, y: S/0.693*2.5
  \draw[very thick,red] plot[smooth,domain=0.02:0.98,samples=60]
    (\x*2.5, {(-(\x*ln(\x)+(1-\x)*ln(1-\x)))/0.6931*2.5});
  \draw[dashed] (1.25,0) -- (1.25,2.5);
  \node[red,above right] at (1.25,2.5) {\scriptsize 최대 $\ln 2$ @ $\theta=\tfrac12$};
  % partial molar arrows (diverging)
  \node at (0.35,1.6) {\scriptsize $\partial S/\partial\theta\to+\infty$};
  \node at (2.15,1.6) {\scriptsize $\to-\infty$};
  \node[below] at (1.25,-0.55) {\small (b) configurational 엔트로피};
\end{scope}
\end{tikzpicture}
\caption{(a) 단일 자리 점유 분포 $\langle n\rangle$~\eqref{eq:occ} — Fermi--Dirac 와 동형이며 Chapter 1 logistic 평형 점유의 기원. (b) configurational 엔트로피 $S_\config(\theta)$~\eqref{eq:Sconfig}; 중심 $\theta=\tfrac12$ 에서 최대 $\ln 2$, 부분몰 $\partial S_\config/\partial\theta=-R\ln[\theta/(1-\theta)]$ 가 양 끝에서 $\pm\infty$ 로 발산(봉우리 내부 $\partial U_\oc/\partial T$ 의 분포 항).}
\label{fig:occ_config}
\end{figure}

\begin{warnbox}
\textbf{이중계산 금지 (파생 B).} 먼저 표기를 못 박는다: \textbf{전이 중심 표준 반응 엔트로피}를
$\Delta S^0_j\equiv\Delta S_{\rxn,j}\equiv[\Delta S(x)]_{\xi=1/2}$ 로 한 기호로 정의한다(봉우리 중심 $\xi=\tfrac12$
에서 측정되는 값, config$=0$). 식~\eqref{eq:dVdT_config} 의 두 항은 \emph{서로 다른 출처}이며 한 번씩만 센다.
첫 항 $\Delta S^0_j$ 는 중심 표준값으로 vibrational$\cdot$electronic 의 중심 기여를 이미 흡수한 상수다. 둘째 항
$+(1/F)\partial S_\config/\partial\theta|_{\theta=1-\xi}=+(R/F)\ln[\xi/(1-\xi)]$ 은 봉우리 \emph{내부}의 $x$-의존을
logistic 폭 $w$ 가 주는 \emph{분포 항}이다. configurational 엔트로피를 $\Delta S^0_j$ 에 \emph{또} 더하면 같은
물리를 두 번 세는 것이다. 측정 $\Delta S(x)$ 의 올바른 분해는
\[
\Delta S(x) \;=\; \underbrace{\Delta S^0_j}_{\text{중심 표준값}}
\;+\;\underbrace{R\ln\frac{\xi}{1-\xi}}_{\text{분포 config (}w\text{ 가 줌)}}
\;+\;\underbrace{\delta S_{\vib/e}(x)}_{\text{전이 폭 내 급변 시만 분리}},
\]
이며 $\Delta S^0_j$ 와 config 항은 결코 겹치지 않는다(중심에서 config$=0$ 이므로 정의가 모순 없이 맞물린다).
여기서 $\delta S_{\vib/e}(x)$ 는 vib$\cdot$electronic 의 \emph{중심값을 뺀} 전이 폭 내 잔차로, 보통 0(중심값에 흡수)
이고 MIT 같은 급변이 있을 때만 살린다(\S\ref{ssec:elec}).
\end{warnbox}

\subsection{문헌 검증 — Chapter 1 staging 엔트로피와의 정합}
이 분포 기반 분해가 실제 흑연 측정과 맞는지 확인한다. 문헌은 흑연 전극의 부분몰 엔트로피 $\Delta S(x)$ 가
저 Li 분율($x<0.2$--$0.25$)에서 양(configurational 우세), 고 분율에서 음(vibrational 우세)이며, stage
$2\to1$($x\approx0.5$)에서 급변을 보인다고 보고한다 \cite{reynier2003,allart2018}. 정량으로는 전극
$\Delta S\approx+29$ J\,mol$^{-1}$K$^{-1}$ @\,$x=0.08$, $-5\sim-16$ J\,mol$^{-1}$K$^{-1}$ @\,$x>0.5$ 이고
\cite{allart2018}, MCMB 흑연 엔트로피 계수가 저-$x$ 에서 $+3\sim4$ mV/K, $\sim0.2$ SOC 부근에서 부호 반전한다
\cite{msmr_partII}. 표~\ref{tab:ds} 는 Chapter 1 의 전이별 중심 표준값 $\Delta S^0_j=+29\to0\to-5\to-16$
J\,mol$^{-1}$K$^{-1}$ 가 이 문헌 프로파일과 부호$\cdot$규모에서 정량 일치함을 보인다 — 곧 분포 spine 의 입력
표준값이 임의값이 아니라 측정에 닿아 있다(양 끝 $+29$@저-$x$, $-16$@$x>0.5$ 정확 대응).

\begin{table}[t]
\centering\small
\renewcommand{\arraystretch}{1.2}
\caption{전이별 중심 표준값 $\Delta S^0_j$ ↔ 문헌 흑연 $\Delta S(x)$(Allart 등 \cite{allart2018}, 전극 분리). 이 표의 $\Delta S^0_j$ 는 봉우리 중심 표준값(파생 B); 봉우리 내부 $x$-의존은 식~\eqref{eq:dSconfig} 의 분포 항이 추가로 줌.}
\label{tab:ds}
\begin{tabular}{l c c c}
\toprule
전이 & $x$ 범위 & $\Delta S^0_j$ [J\,mol$^{-1}$K$^{-1}$] & 문헌 $\Delta S(x)$ \\
\midrule
$4\!\to\!3$ & 0.08–0.16 & $+29.0$ & $+29$ @\,$x{=}0.08$ (양 끝, 정확) \\
$3\!\to\!2\mathrm L$ & 0.16–0.25 & $0.0$ & 중간 하강 구간 $0\!\to\!-16$ (범위) \\
$2\mathrm L\!\to\!2$ & 0.25–0.50 & $-5.0$ & 중간 하강 구간 (범위) \\
$2\!\to\!1$ & 0.50–1.00 & $-16.0$ & $-16$ @\,$x{>}0.5$ (양 끝, 정확) \\
\bottomrule
\end{tabular}

\smallskip
{\footnotesize $\dagger$ Chapter 1 의 4 주요 plateau 전이와 Allart 등 \cite{allart2018} 의 세분 region(dilute
LiC$_{24}$ 포함, 양 끝 reg IV/reg I)은 $x$-경계가 $1{:}1$ 이 아니다. \emph{양 끝}($+29$ @저-$x$, $-16$ @$x{>}0.5$)
은 정확 대응하고, 중간 두 전이는 $\Delta S$ 가 $0\to-16$ 으로 하강하는 \emph{구간 안}에 놓인다(점-대-점이 아닌
범위 정합). 해상도 차이일 뿐 추세$\cdot$부호는 일관.}
\end{table}

중심 표준값 $\Delta S^0_j$ 는 봉우리 중심에서 측정되는 상수이고, 부분몰 config 항은 그 위에 분포가 더하는
$x$-의존이라는 점을 다시 강조한다(파생 B). 한편 엔탈피 쪽 표준값 $\Delta H^0_j$ 는 \emph{기준 차이}에 주의해야
한다: Chapter 1 의 $\Delta H_{\rxn,j}$ 는 \emph{전이별}(differential) 반응 엔탈피이고, calorimetry 의 형성
엔탈피(LiC$_6$ $-13.9$, LiC$_{12}$ $-24.8$ kJ/mol\,Li \cite{chemmater2015})는 \emph{누적}(formation,
graphite+Li 금속 기준)이라 직접 등치할 수 없다 — 비교하려면 형성 엔탈피의 $x$-미분으로 환산해야 한다. 전이
중심에서는 $\Delta H^0_j=-FU_j+FT\,(\partial U_j/\partial T)$ 가 표준값들끼리 일관된다(예: stage $2\to1$ 에서
$-13.0$ kJ/mol, 표값 일치).

% ====================================================================
\section{vibrational(포논)·electronic(전자) 엔트로피 — 두 분포의 추가}\label{sec:vibel}
% ====================================================================
configurational 엔트로피는 Li 이온의 \emph{배열} 분포였다. 그러나 격자에는 또 다른 두 자유도가 있고, 각각
자신의 분포를 갖는다 — 격자 진동(포논)은 Bose--Einstein 분포를, 전도 전자는 Fermi--Dirac 분포를 따른다.
이 절은 같은 분포 언어로 두 항을 세워, 세 분포가 합쳐져 한 챕터의 엔트로피를 이룸을 보인다.

\subsection{vibrational 엔트로피 — 포논 Bose--Einstein 분포}
삽입은 격자 상수와 결합 강도를 바꾸어 포논 스펙트럼을 이동시킨다. 진동수 $\omega_k$ 의 한 조화 모드는 보손이고,
평균 점유는 Bose--Einstein 분포 $n_k=1/(e^{\beta\hbar\omega_k}-1)$ 이다. 한 모드의 엔트로피는 보손 통계에서
\begin{equation}
S_{\vib,k} \;=\; \kB\bigl[(1+n_k)\ln(1+n_k) - n_k\ln n_k\bigr],
\label{eq:Svib_mode}
\end{equation}
이고, 전 모드 합(1몰당, $R=N_A\kB$)이 $S_\vib=R\sum_k[(1+n_k)\ln(1+n_k)-n_k\ln n_k]$ 다. 삽입에 따른 부분몰
변화 $\Delta S_\vib=\partial S_\vib/\partial x$ 는 모드 연화/경화에 따라 부호가 갈리는데, 흑연에서는 고-$x$
(만충에 가까운) 영역에서 \emph{음의 baseline}을 만든다(Reynier 의 ``second contribution'') \cite{reynier2003}.
저-$x$ 에서는 configurational 양의 발산이 지배하고, 고-$x$ 에서 이 vibrational 음 baseline 이 드러난다. 제일원리
계산도 흑연 삽입 엔트로피를 vibrational$\cdot$configurational 기여로 분해해 같은 추세를 보고한다 \cite{jpcc2021}.

전이 폭 안에서 $\Delta S_\vib$ 가 천천히 변하면(흑연의 통상적 경우) 그 값은 봉우리 \emph{중심 표준값}
$\Delta S^0_j$ 에 흡수된다 — 곧 표~\ref{tab:ds} 의 고-$x$ 음수값($-5,-16$ J\,mol$^{-1}$K$^{-1}$)이 이미
vibrational 기여를 품고 있다. 이것이 파생 B 에서 ``$\Delta S^0_j$ 가 vib$\cdot$electronic 의 중심값을 흡수''라고
말한 의미다. config 항만 분포가 따로 떼어 $x$-의존으로 살린다.

\subsection{electronic 엔트로피 — Fermi--Dirac 분포와 Sommerfeld 항}\label{ssec:elec}
금속성 호스트에서는 전도 전자도 엔트로피를 진다. Fermi--Dirac 분포 $f(E)=1/(e^{\beta(E-E_F)}+1)$ 에서, 저온
($\kB T\ll E_F$) Sommerfeld 전개로 전자 엔트로피는
\begin{equation}
S_e \;=\; \frac{\pi^2}{3}\,\kB^2\,T\,g(E_F),
\label{eq:Se}
\end{equation}
가 되어 온도에 \emph{선형}이고 Fermi 준위 상태밀도 $g(E_F)$ 에 비례한다. 부분몰 전자 엔트로피
$\Delta S_e=\partial S_e/\partial x$ 는 삽입이 $g(E_F)$ 와 $E_F$ 를 바꾸는 방식에 달려 있다. Chapter 1 v9 의
규약대로, 삽입(리튬화)은 일반적으로 $\Delta S_e<0$ 이다 — 밴드 채움으로 $E_F$ 부근 상태밀도가 줄어드는 호스트가
많기 때문이다.
\begin{itemize}[leftmargin=1.4em,itemsep=1pt]
\item \textbf{흑연 하프셀}: 흑연은 준금속이라 $g(E_F)$ 가 작아 $\Delta S_e\approx0$ — 전자 항은 소수이고
config$+$vib 가 지배한다. 본 장 주 사례에서 electronic 은 보정항이다.
\item \textbf{LCO 하프셀(사례)}: Li$_x$CoO$_2$ 는 $x$ 에 따라 금속--절연체 전이(MIT)를 겪어 $g(E_F)$ 가
급변한다. 그 전이 폭 안에서 $\Delta S_e$ 가 \emph{급변}하므로 중심 표준값에 흡수되지 못하고 $x$-의존으로
유지된다(파생 I 의 예외, \S\ref{sec:limits}). 이것이 electronic 항을 분포로 다뤄야 하는 이유의 본보기다.
\end{itemize}

\begin{keybox}
세 분포가 한 챕터의 엔트로피를 이룬다: Li 배열의 \emph{격자기체}(config), 격자 진동의 \emph{Bose--Einstein}
(vib), 전도 전자의 \emph{Fermi--Dirac}(electronic). 측정 부분몰 엔트로피는
\[
\Delta S(x)=\underbrace{R\ln\frac{\xi}{1-\xi}}_{\text{config 분포}}
+\underbrace{\Delta S^0_j}_{\text{vib+elec 중심 흡수}}
+\underbrace{\delta S_{\vib/e}(x)}_{\text{전이 폭 내 급변 시만}},
\]
이고, 가역 발열은 이 세 분포를 재배열하는 열이다. config 는 봉우리 \emph{내부}를 비선형으로 만들고
($\pm\infty$ 발산), vib 는 고-$x$ 음 baseline 을, electronic 은 (MIT 가 없으면) 작은 상수를 준다.
\end{keybox}

\subsection{반응 엔트로피 vs 활성화 엔트로피 — 혼동 금지}
가역 발열에 들어가는 것은 위 세 분포가 주는 \emph{반응(평형) 엔트로피} $\Delta S(x)$ 뿐이다. 이는 Chapter 1
동역학 꼬리의 \emph{활성화 엔트로피} $\Delta S_{a,j}$(Eyring 속도 $k_j\propto\exp[-(\Delta H_a-T\Delta S_a)/RT]$
의 prefactor 온도의존)와 \emph{차원은 같으나 물리가 다른 양}이다. 활성화 엔트로피는 \emph{비가역} 동역학을
지배하고 가역 발열에는 들어가지 않는다 \cite{msmr_partI}. 혼동하면 가역열을 동역학 lag 와 섞게 된다.

% ====================================================================
\section{섞임과 겹침 — 분포가 만드는 $\partial U_\oc/\partial T(x)$ 의 모양}\label{sec:mixing}
% ====================================================================
실제 흑연 곡선은 단일 전이가 아니라 여러 전이가 \emph{겹쳐} 있다. 한 SOC 에서 여러 전이의 점유 분포가 동시에
변하고 있고, 측정 $\partial U_\oc/\partial T(x)$ 는 그 분포들의 가중 합으로 나타난다. 이 절은 (A) 겹침 가중식을
음함수 미분으로 닫고 수치로 검증하며, (B) 이중계산을 분리하고, (C) 상호작용 $\Omega$ 가 봉우리 폭과 엔트로피
모양을 어떻게 변형하는지, (D) 히스테리시스 분기별 $\partial U/\partial T$ 를 가역/비가역으로 어떻게 가르는지를
닫힌식으로 다룬다.

\subsection{파생 A — 겹침의 $dQ/dV$-비중 가중 (계단이 아닌 연속 블렌드)}\label{ssec:overlap}
관측 평형 전위 $U_\oc(x,T)$ 는 모든 전이의 점유가 합쳐 전하 보존을 만족하는 음함수로 정의된다:
\begin{equation}
\sum_j Q_j\,\xi_{\eq,j}(U_\oc,T) \;=\; Q\,x,
\label{eq:implicit}
\end{equation}
여기서 $Q_j$ 는 전이 $j$ 의 용량, $Q=\sum_j Q_j$ 다. 양변을 고정 $x$ 에서 $T$ 로 음함수 미분하면
\begin{equation}
\frac{\partial U_\oc}{\partial T}\Big|_x
\;=\; -\,\frac{\sum_j Q_j\,\partial\xi_j/\partial T|_U}{\sum_j Q_j\,\partial\xi_j/\partial U|_T}.
\label{eq:implicit_diff}
\end{equation}
logistic 점유~\eqref{eq:logistic} 에서 두 편미분이 닫힌형으로 나온다. 전위 미분은 곧 \emph{국소 $dQ/dV$ 봉우리
모양}이다:
\begin{equation}
\frac{\partial\xi_j}{\partial U}\Big|_T \;=\; \frac{\xi_j(1-\xi_j)}{w_j} \;\equiv\; g_j(x),
\label{eq:gj}
\end{equation}
온도 미분은 \emph{두 조각}을 갖는다 — 봉우리 중심 $U_j(T)$ 의 온도 이동과, 폭 $w(T)=RT/F$ 의 명시적 온도
의존이다. $\xi_j$ 가 logistic 인자 $a_j=(U-U_j(T))/w(T)$ 의 함수이고 $\partial\xi_j/\partial a_j=g_j w_j$,
$\partial w/\partial T=R/F$ 임을 쓰면($z_j\equiv(U-U_j)/w_j=\ln[\xi_j/(1-\xi_j)]$)
\begin{equation}
\frac{\partial\xi_j}{\partial T}\Big|_U
\;=\; -\,g_j(x)\,\frac{\partial U_j}{\partial T}
\;\underbrace{-\,g_j(x)\,\frac{R}{F}\,z_j}_{w(T)\text{ 조각}},
\qquad z_j=\ln\frac{\xi_j}{1-\xi_j}.
\label{eq:dxidT}
\end{equation}
첫 조각만(중심 이동)을 \eqref{eq:gj}$\cdot$\eqref{eq:dxidT} 와 함께 \eqref{eq:implicit_diff} 에 넣으면, 분모의
$\partial\xi_j/\partial U=g_j$ 와 분자의 $-g_j\partial U_j/\partial T$ 가 만나 겹침 가중의 \emph{중심값 식(단순식)}이
닫힌다:
\begin{equation}
\boxed{\;\frac{\partial U_\oc}{\partial T}(x)\Big|_{\text{단순식}}
\;=\; \frac{\sum_j Q_j\,g_j(x)\,(\partial U_j/\partial T)}{\sum_j Q_j\,g_j(x)}
\;=\; \frac1F\,\frac{\sum_j Q_j\,g_j(x)\,\Delta S_{\rxn,j}}{\sum_j Q_j\,g_j(x)}\;.}
\label{eq:weighted}
\end{equation}
곧 단순식의 관측 엔트로피 계수는 각 전이의 $\Delta S_{\rxn,j}/F$ 를 \emph{국소 $dQ/dV$ 비중}
$Q_jg_j/\sum_k Q_kg_k$ 로 가중한 평균이다. 겹침 영역에서 인접한 $g_j,g_{j+1}$ 둘 다 0 이 아니므로, $x$ 가 한
전이에서 다음으로 넘어갈 때 $\Delta S_{\rxn,j}/F$ 에서 $\Delta S_{\rxn,j+1}/F$ 로 \emph{연속}으로 이동한다 —
\textbf{계단이 아니라 연속 블렌드}다(그림~\ref{fig:blend}). 가중 분모 $\sum Q_jg_j$ 가 정확히 측정 $dQ/dV$(=$\partial x/\partial U$)
이므로, 비중은 ``그 SOC 에서 어느 봉우리가 활성인가''를 그대로 읽는다. \eqref{eq:dxidT} 의 둘째 조각($w(T)$
의존, $-g_jRz_j/F$)을 마저 넣으면 봉우리 내부 config 항 $+R z_j/F=+R\ln[\xi_j/(1-\xi_j)]/F$ 가 더해진
\emph{완전식}이 되며(\S\ref{sec:config} 의 부분몰 config 와 동일), 이것이 다음 절 파생 B 와 수치 검증의 대상이다.

\begin{srcbox}
★\textbf{수치 검증(파생 A).} Chapter 1 코드(\texttt{Anode\_Fit\_v11})의 4-전이 흑연 staging 파라미터로,
유한차분 $\partial U_\oc/\partial T|_x^\mathrm{FD}=[U_\oc(x,T_2)-U_\oc(x,T_1)]/\Delta T$ ($T_1{=}292.15$,
$T_2{=}298.15$ K)를 식~\eqref{eq:weighted} 의 \emph{단순식}(중심값 $\Delta S^0_j/F$ 만)과, 그리고 봉우리 내부
config 항 $z_j R/F$ ($z_j=(U_\oc-U_j)/w_j$)를 포함한 \emph{완전식}과 비교했다 \cite{numverif2026}. 결과:
완전식은 175 점 전 범위에서 FD 와 \emph{부동소수점 정밀도}로 일치(절대오차 $\approx0$ mV/K), 단순식은 봉우리
내부 config 항만큼 체계적으로 빗나간다(단순식 절대오차 최대 $0.18$ mV/K). 그 config 항 자체의 부호 있는 크기는
구간에 따라 $[-0.21,+0.14]$ mV/K 범위다. 전이 경계 4 곳 모두 \emph{계단 없이} 연속 블렌드(인접 $\Delta S^0/F$
사이 값을 취함)임도 확인됐다. 곧 \eqref{eq:weighted} 의 중심값 가중에 \S\ref{sec:config} 의
분포 config 를 더하면, 임시 계단 함수나 외부 입력 없이 \emph{측정급 비선형} $\partial U_\oc/\partial T(x)$ 가
모델 안에서 자동 생성된다.
\end{srcbox}

\begin{figure}[t]
\centering
\begin{tikzpicture}[x=1.0cm,y=1.0cm]
  % axes
  \draw[->,>=Stealth] (-0.2,0) -- (9.4,0) node[right]{\small $x$ (탈리튬화 진행)};
  \draw[->,>=Stealth] (0,-2.0) -- (0,2.3) node[above]{\small $\partial U_\oc/\partial T$};
  \draw[dashed,gray] (0,0) -- (9.2,0);
  % two transition levels dS_j/F (arbitrary units): j has +1.5, j+1 has -1.5
  \draw[densely dotted] (0.4,1.5) -- (3.0,1.5) node[right,black]{\scriptsize $\Delta S^0_{j}/F$};
  \draw[densely dotted] (6.0,-1.5) -- (8.8,-1.5) node[right,black]{\scriptsize $\Delta S^0_{j+1}/F$};
  % continuous blend curve: smooth sigmoid from +1.5 to -1.5 centered at x=4.6
  \draw[very thick,blue] plot[smooth,domain=0.4:8.8,samples=80]
    (\x, {1.5 - 3.0/(1+exp(-1.6*(\x-4.6)))});
  \node[blue,above] at (2.0,1.6) {};
  % the overlap region shading
  \draw[decorate,decoration={brace,amplitude=4pt},gray] (3.2,-1.9) -- (5.9,-1.9);
  \node[gray,below] at (4.55,-2.05) {\scriptsize 겹침 영역: 연속 블렌드};
  % weight ratio annotation
  \node[blue] at (6.7,1.0) {\scriptsize $\dfrac{\sum_j Q_jg_j\,\Delta S^0_j/F}{\sum_j Q_jg_j}$};
  % NOT a step: contrast dashed step
  \draw[red,dashed,thick] (0.4,1.5) -- (4.6,1.5) -- (4.6,-1.5) -- (8.8,-1.5);
  \node[red] at (3.7,-0.55) {\scriptsize 계단(틀림)};
\end{tikzpicture}
\caption{겹침 가중식~\eqref{eq:weighted} 이 만드는 \emph{연속 블렌드}(파란 실선): $x$ 가 전이 $j$ 에서 $j{+}1$ 로 넘어갈 때 $\partial U_\oc/\partial T$ 가 $\Delta S^0_j/F$ 에서 $\Delta S^0_{j+1}/F$ 로 국소 $dQ/dV$ 비중 $Q_jg_j/\sum Q_kg_k$ 에 따라 \emph{연속}으로 이동한다. 계단(빨간 점선)이 아니다 — 수치 검증에서 전이 경계 4 곳 모두 연속임이 확인됐다 \cite{numverif2026}.}
\label{fig:blend}
\end{figure}

\subsection{파생 B — 봉우리 내부 configurational 과 이중계산 분리}
완전식이 FD 와 정확히 일치하는 이유는, \S\ref{sec:config} 에서 유도한 부분몰 config 항이 정확히 빠진
``단순식의 잔차''이기 때문이다. 단일 전이가 지배하는 영역에서 식~\eqref{eq:dVdT_config} 가 그대로 살아난다:
\begin{equation}
\frac{\partial U_\oc}{\partial T}(x)\;\approx\;\frac{\Delta S_{\rxn,j}}{F}\;+\;\frac{R}{F}\ln\!\frac{\xi_j}{1-\xi_j}
\qquad(\text{단일 전이 지배}).
\label{eq:single_config}
\end{equation}
여기서 두 항이 별개 출처임을 다시 확인한다(파생 B): $\Delta S^0_j/F$ 는 봉우리 \emph{중심}($\xi_j=\tfrac12$,
config$=0$)의 표준값, $R\ln[\xi_j/(1-\xi_j)]/F$ 는 봉우리 \emph{내부}의 분포 항이다. 수치 검증에서 ``완전식
$-$ 단순식'' 차이가 바로 이 config 항이었고(부호 있는 크기 $[-0.21,+0.14]$ mV/K), 이는 \S\ref{sec:config} 의
$+(1/F)\partial S_\config/\partial\theta|_{\theta=1-\xi}$ 와 같은 식이다. \textbf{config 를 $\Delta S^0_j$ 에 또 더하면
이중계산}이며, 올바른 분해는 ``중심 표준값 $+$ 분포 config''로 정확히 한 번씩 센다. 봉우리 중심에서 config$=0$
이므로 두 정의가 모순 없이 맞물린다(중심에서 측정한 표준값에 config 가 겹쳐 들어가지 않는다).

\subsection{파생 C — 상호작용에 의한 유효 폭 $w_\eff(\Omega)$}\label{ssec:weff}
이상 격자기체는 봉우리 폭이 $w=RT/F$ 로 고정이지만, regular-solution 상호작용 $\Omega$ 은 폭을 바꾼다.
식~\eqref{eq:slope_BW} 의 평형 등온선 기울기를 봉우리 중심 $\theta=\tfrac12$ 에서 보면
\begin{equation}
\frac{\dd V_\eq}{\dd\theta}\Big|_{1/2}=\frac{4RT-2\Omega}{F}
\;=\;\frac{4}{F}\Bigl(RT-\tfrac{\Omega}{2}\Bigr),
\label{eq:slope_half}
\end{equation}
이상($\Omega=0$)의 중심 기울기 $4RT/F$ 와 비교하면, 상호작용이 기울기를 인자 $(1-\Omega/2RT)$ 로 줄인다.
이상 logistic 의 중심 기울기는 $\dd V/\dd\xi|_{1/2}=4w$ 로 \emph{폭에 비례}하므로($w=RT/F$), 더 작은
기울기를 같은 logistic 형으로 재현하려면 \emph{유효 폭이 더 작아야} 한다. 곧 중심 기울기를 정합시키는
유효 폭은
\begin{equation}
\boxed{\;w_\eff(\Omega)\;=\;w\Bigl(1-\frac{\Omega}{2RT}\Bigr)\;=\;\frac{RT}{F}\Bigl(1-\frac{\Omega}{2RT}\Bigr)\qquad(0\le\Omega<2RT)\;.}
\label{eq:weff}
\end{equation}
$\Omega>0$(인력, 상분리 경향)이면 $w_\eff<w$ — 봉우리가 \emph{좁아진다}(그림~\ref{fig:weff}; 중심 기울기 $\dd V_\eq/\dd\theta$ 가
작아져 같은 용량을 더 좁은 전위 폭에 쏟고, $dQ/dV$ 봉우리는 더 뾰족$\cdot$높아진다). 이 좁아진 봉우리는 부분몰
config 엔트로피 모양도 변형한다: 같은 전위 폭 안에 $\ln[\xi/(1-\xi)]$ 가 더 가파르게 변하므로 $\partial U_\oc/\partial T$
의 \emph{봉우리 내부 변동이 증폭}된다. $\Omega\to2RT^-$ 에서 $w_\eff\to0$, 중심 기울기 $\to0$ — 이것이 \emph{상분리
임계}이고, $dQ/dV$ 봉우리가 한 전위로 무한히 좁아진(delta 형) plateau 가 된다. 그 너머($\Omega>2RT$)는 전위
평탄$\cdot$불연속과 $\partial U_\oc/\partial T$ 의 평탄$\cdot$불연속을 낳는다(\S\ref{sec:limits}).
\begin{warnbox}
\textbf{$\Omega$ 의 온도 의존(2차 항).} 위는 $\Omega$ 가 $T$-무관일 때다. $\Omega$ 자체가 약한 온도 의존
($\partial\Omega/\partial T\ne0$)을 가지면, $\partial w_\eff/\partial T$ 가 $\partial U_\oc/\partial T$ 에 2차
기여를 더한다. 흑연에서는 보통 소수이나, 정밀 분석에서는 명시해야 한다(상호작용 엔트로피 $\propto\partial\Omega/\partial T$).
\end{warnbox}

\begin{figure}[t]
\centering
\begin{tikzpicture}[x=1.0cm,y=1.0cm]
  \draw[->,>=Stealth] (-0.2,0) -- (6.0,0) node[right]{\small $\Omega/(2RT)$};
  \draw[->,>=Stealth] (0,-0.2) -- (0,3.6) node[above]{\small $w_\mathrm{eff}/w$};
  % gridline at omega/2RT = 1 (critical)
  \draw[dashed,red] (5.0,0) -- (5.0,3.2); \node[red,below right] at (5.0,-0.05) {\scriptsize $1$ (임계)};
  \draw[dashed,gray] (0,3.0) -- (5.6,3.0); \node[left] at (-0.05,3.0){\scriptsize $1$};
  % w_eff/w = 1-u, u in [0,1], x: u*5, y: value*3 (so 1->height3, 0->0)
  \draw[very thick,blue] plot[smooth,domain=0:1.0,samples=40]
    (\x*5, {(1-\x)*3});
  \node[blue] at (1.9,1.55) {\scriptsize $w_\mathrm{eff}=w\bigl(1-\tfrac{\Omega}{2RT}\bigr)$};
  % critical point dot at (5,0)
  \fill[red] (5.0,0) circle (1.6pt);
  \node[red,above right] at (5.0,0.05) {\scriptsize $w_\mathrm{eff}\to0$};
  % regimes
  \node[align=center] at (1.3,-0.7) {\scriptsize 균질 고용체\\\scriptsize $\Omega<2RT$};
  \node[align=center,red] at (5.5,1.2) {\scriptsize 상분리\\\scriptsize ($\Omega>2RT$)};
\end{tikzpicture}
\caption{유효 봉우리 폭 $w_\eff(\Omega)$~\eqref{eq:weff}: 인력 상호작용($\Omega>0$)이 중심 기울기를 줄여 봉우리를 \emph{좁힌다}($w_\eff<w$ — 같은 용량을 더 좁은 전위 폭에, $dQ/dV$ 봉우리는 더 뾰족), 봉우리 내부 config 엔트로피 변동을 증폭한다. $\Omega\to2RT^-$ 에서 $w_\eff\to0$, 중심 기울기 $\to0$ — 상분리 임계($dQ/dV$ 가 한 전위로 무한히 좁아진 plateau; 그 너머는 $\partial U_\oc/\partial T$ 의 평탄$\cdot$불연속).}
\label{fig:weff}
\end{figure}

\subsection{파생 D — 히스테리시스: 충/방전 분기와 가역/비가역 분리}\label{ssec:hys}
흑연 OCV 는 충방전 사이 경로의존이다(준안정 stacking, stage II 무질서) \cite{hysteresis2018}. 분기별 전이 중심을
$U_j^{(d)}=U_j\pm\tfrac12\Delta U_j^\hys$ (방향 $d=$ ch/dis)로 두면, 각 분기의 엔트로피 계수는 \emph{그 분기의}
봉우리 모양 $g_j^{(d)}$ 로 가중된 \eqref{eq:weighted} 다:
\begin{equation}
\frac{\partial U_\oc^{(d)}}{\partial T}
\;=\;\frac1F\,\frac{\sum_j Q_j\,g_j^{(d)}(x)\,\Delta S_{\rxn,j}}{\sum_j Q_j\,g_j^{(d)}(x)}.
\label{eq:hys_branch}
\end{equation}
가역 발열에 들어가는 것은 두 분기의 \emph{평균}이다 — 충방전을 한 사이클 돌면 히스 중심이 상쇄되고 열역학적
$\Delta S$ 만 남기 때문이다:
\begin{equation}
\boxed{\;\frac{\partial U_\oc^\rev}{\partial T}
\;=\;\tfrac12\Bigl(\frac{\partial U_\oc^\mathrm{ch}}{\partial T}+\frac{\partial U_\oc^\mathrm{dis}}{\partial T}\Bigr)\;.}
\label{eq:hys_rev}
\end{equation}
히스 gap $\Delta U^\hys$ 자체는 \emph{비가역} 과정으로, 한 사이클당 $\propto I\,\Delta U^\hys$ 의 entropy
production(소산열)을 만든다 — 이는 온도 미분 $\partial/\partial T$ 와 \emph{별개}이며 가역열에 섞이지 않는다.
\begin{keybox}
\textbf{혼동 금지.} 가역 발열은 분기 평균~\eqref{eq:hys_rev} 로, 히스 gap 은 비가역열로 \emph{분리} 계산한다.
$\partial U^\rev/\partial T$ 는 분포 재배열의 가역 엔트로피이고, $\Delta U^\hys$ 는 경로 비가역성의 소산이다.
하나의 $\partial U/\partial T$ 로 둘을 뭉뚱그리면 가역/비가역 발열을 섞게 된다(경로의존 측정 불확실도의
정량은 본 장 범위 밖, 공백으로 명시 \cite{hysteresis2018}).
\end{keybox}

% ====================================================================
\section{극한과 코너 — 분포 검산과 ``$\Delta S_{\rxn,j}$ 상수'' 근사의 타당성}\label{sec:limits}
% ====================================================================
닫힌식의 신뢰는 극한에서 옳은 물리로 환원되는지로 검산된다. 표~\ref{tab:limits} 는 여섯 코너를 정리한다 —
각각 분포 식이 물리적으로 맞는 값으로 가는지 확인한다.

\begin{table}[t]
\centering\small
\renewcommand{\arraystretch}{1.35}
\caption{$\partial U_\oc/\partial T(x)$ 의 극한$\cdot$코너 검산. 규약: $\xi=$ 탈리튬화 진행률($\xi\to1$ 희박, $\xi\to0$ 만충).}
\label{tab:limits}
\begin{tabular}{l l l}
\toprule
극한 & $\partial U_\oc/\partial T$ 거동 & 물리$\cdot$검증 \\
\midrule
$\xi\to1$ (희박, Li-poor) & config $\to+\infty$ & 삽입 자리 선택지 폭증; 저-$x$ 큰 양 $\Delta S$($+29$)의 주원천 \\
$\xi\to0$ (만충, Li-rich) & config $\to-\infty$ & 배열 선택지 소멸; 만충 정렬의 음 부분몰 엔트로피 \\
$\xi=\tfrac12$ (봉우리 중심) & $=\Delta S_{\rxn,j}/F$ & config$=0$; 중심 표준값 정확 회수(파생 B 일관) \\
$\Omega\to2RT^-$ & $w_\eff\to0$, plateau & 상분리 임계; $dQ/dV$ delta화, $\partial U_\oc/\partial T$ 평탄$\cdot$불연속 \\
단일 봉우리 (겹침 0) & $=\Delta S_{\rxn,j}/F+$ config & 가중식~\eqref{eq:weighted} 이 단일 전이로 환원 \\
고온 ($\kB T\!\sim\!E_F$ 근접) & electronic $\propto T$ 우세화 & FD Sommerfeld~\eqref{eq:Se}; $S_e$ 선형 증가 \\
\bottomrule
\end{tabular}
\end{table}

여섯 코너 모두 분포 식이 옳은 극한으로 환원된다. 특히 세 가지가 본 장의 골격을 닫는다: (i) 중심
$\xi=\tfrac12$ 에서 config$=0$ 이라 표준값이 정확히 회수되므로 파생 B 의 분리가 자기일관적이다; (ii) 단일
봉우리 한계에서 가중식이 단일 전이 $+$ config 로 환원되므로 \eqref{eq:weighted} 가 \eqref{eq:single_config} 를
포함한다; (iii) $\Omega\to2RT$ 에서 $w_\eff\to0$(봉우리가 한 전위로 무한히 좁아짐)이 상분리 plateau 와 맞물려,
상호작용 모형이 상전이 임계를 스스로 짚는다.

\begin{keybox}
★\textbf{``$\Delta S_{\rxn,j}$ 전이당 상수'' 근사의 판정.} 측정 $\partial U_\oc/\partial T(x)$ 의 비선형은 두
출처에서 \emph{자동 생성}된다 — (i) 겹침 가중(\S\ref{ssec:overlap})과 (ii) 봉우리 내부 config 분포
(\S\ref{sec:config}). 따라서 $\Delta S_{\rxn,j}(x)$ 를 임의 함수로 둘 필요가 \emph{없다}: 전이당 \emph{상수}
표준값 $+$ 분포만으로 측정급 곡선이 나온다. 이 근사는 \emph{중심 표준값(vib$+$electronic 중심)이 전이 폭 안에서
천천히 변할 때} 타당하다. 예외는 그 항이 전이 폭 \emph{안에서} 급변하는 경우 — 대표적으로 LCO 의 MIT 가 한
전이 안에서 $g(E_F)$ 를 급변시키면 $\Delta S_e(x)$ 만은 $x$-의존으로 유지해야 한다(\S\ref{ssec:elec}). 흑연
하프셀에서는 electronic 이 소수라 이 예외가 거의 없고, config$+$중심 상수로 충분하다.
\end{keybox}

% ====================================================================
\section{가역 발열 — 분포를 재배열하는 열}\label{sec:revheat}
% ====================================================================
이제 분포에서 세운 $\partial U_\oc/\partial T(x)$ 를 가역 발열로 닫는다. 전지 셀의 발열은 Bernardi--Pawlikowski--
Newman 의 일반 에너지 수지 \cite{bernardi1985} 에서 가역(엔트로피)과 비가역(소산) 두 성분으로 갈린다. 단일
활물질 한계에서
\begin{equation}
\dot Q \;=\; \underbrace{I\,(U_\oc-V)}_{\dot Q_\irr\ge0}\;-\;\underbrace{I\,T\,\frac{\partial U_\oc}{\partial T}}_{\dot Q_\rev},
\qquad
\boxed{\;\dot Q_\rev \;=\; -\,I\,T\,\frac{\partial U_\oc}{\partial T}\;=\;-\,\frac{I\,T}{F}\,\Delta S(x)\;,}
\label{eq:qrev}
\end{equation}
이다($I>0$ 방전, $V$ 단자 전압, $U_\oc$ 평형 전위). 첫 항은 과전압 소산(2법칙상 항상 발열, Chapter 1 의
동역학 꼬리$\cdot$분극이 만드는 entropy production), 둘째 항이 가역열로 부호 양방향이다.
\begin{srcbox}
부호 규약: 평형 전위와 반응 깁스 에너지가 $\Delta G=-FU_\oc$ 로 묶이고 $\Delta S=-\partial(\Delta G)/\partial T
=+F\,\partial U_\oc/\partial T$ 이므로 $\dot Q_\rev=-IT\,\partial U/\partial T=-(IT/F)\,T\,\Delta S/T=-(IT/F)\Delta S$
다 \cite{newman,msmr_partI}. 본 장은 평형 전위 $U_\oc$ 기준 $\Delta S=+F\,\partial U_\oc/\partial T$ 규약으로 통일한다.
\end{srcbox}

식~\eqref{eq:qrev} 의 $\Delta S(x)$ 는 본 장이 세운 세 분포의 합이다(\S\ref{sec:vibel} 핵심 박스). 곧 가역
발열은 한 충전상태에서 다른 충전상태로 넘어갈 때 \emph{Li 배열 분포$\cdot$포논 분포$\cdot$전자 분포를 재배열하며
주고받는 열}이다. 부호는 식~\eqref{eq:qrev} 가 곧장 정한다 — 방전($I>0$)에서 $\dot Q_\rev=-(IT/F)\Delta S$
이므로 $\Delta S>0$ 이면 $\dot Q_\rev<0$ 으로 \emph{흡열}(endothermic, 셀이 열을 흡수), $\Delta S<0$ 이면
$\dot Q_\rev>0$ 으로 \emph{발열}(exothermic)이다($\Delta S$ 는 삽입 $n=1$ 기준 Li 1몰당
J\,mol$^{-1}$K$^{-1}$). SOC 에 따라 $\Delta S(x)$ 의 부호가 바뀌므로 흡$\cdot$발열이 교대한다: 저-$x$ 에서
config 양의 발산이 지배해 $\Delta S>0$ 이므로 방전 시 $\dot Q_\rev<0$(흡열), 고-$x$ 에서 vibrational 음
baseline 이 드러나 $\Delta S<0$ 이 되어 방전이 발열로 돈다. stage $2\to1$($x\approx0.5$)의 큰 음의 $\Delta S$
는 방전 시 $\dot Q_\rev>0$, 곧 가역 발열의 두드러진 신호를 남기며, 이는 calorimetry 의 가역 발열 직접 관측과
정합한다 \cite{standardised2024}. 충/방전 대칭(가역)이고, 비선형 $\partial U_\oc/\partial T(x)$ 는 분포의
$x$-의존을 직접 반영한다 — 모양이 곧 분포다.

\begin{keybox}
\textbf{계산용 종합식(use-this).} 하프셀 가역 발열을 실제로 산출할 때는 겹침 가중(단순식)에 봉우리 내부 분포
config 를 더한 \emph{완전식}을 쓴다:
\[
\boxed{\;\frac{\partial U_\oc}{\partial T}(x)
=\frac{\sum_j Q_j\,g_j(x)\,\bigl[\Delta S^0_j/F+(R/F)\ln(\xi_j/(1-\xi_j))\bigr]}{\sum_j Q_j\,g_j(x)},\quad
g_j=\frac{\xi_j(1-\xi_j)}{w_j}\;}
\]
$\dot Q_\rev=-IT\,\partial U_\oc/\partial T$. \emph{입력} = $\{\Delta S^0_j,\,Q_j,\,U_j,\,w_j\}$ 와 Ch1 의
$\xi_j(V,T)$(eq~\eqref{eq:logistic}). $\Delta S^0_j$ 는 다온도 $dQ/dV$ 동시 피팅으로 얻고, 나머지는 Ch1 모수다.
단일 전이 지배 영역에서 이 식은 eq~\eqref{eq:single_config} 로 환원된다.
\end{keybox}

\textbf{한계$\cdot$갭(정직).} 본 장의 닫힌식에는 다음 공백이 남는다. (1) 봉우리 내부 config 항은 단일 전이
지배 영역에서 수치로 검증됐으나(파생 A), 강한 겹침 영역의 더 높은 차수 보정은 실데이터 round-trip 으로
확정해야 한다(현재는 선두차수까지다). (2) 히스테리시스 하 경로의존 $\partial U_\oc/\partial T$ 의 측정 불확실도
정량화는 본 장 범위 밖이며, 파생 D 는 가역/비가역 분리의 틀만 제시한다 \cite{hysteresis2018}. (3) 상호작용
$\Omega$ 의 온도 의존(파생 C 의 2차 항)과 그 흑연에서의 값은 아직 확보하지 못했다. (4) electronic 항의 절대값은
흑연에서 소수로 취급했으며, LCO 의 MIT 예외는 분포 언어의 사례로만 다뤘다. (5) 전셀 합성은 본 장의 범위
밖이고, 그 근거는 서두 범위 박스에 명시했다.

\begin{procedurebox}
\textbf{방법론 출구 — 다온도 $dQ/dV$ 에서 $\partial U_j^0/\partial T$ 추정 절차.} 위 종합식의 입력
$\{\Delta S^0_j,Q_j,U_j,w_j\}$ 중 전이별 중심 표준값 $\Delta S^0_j$ 는 다음 순서로 추정한다.
\begin{enumerate}[leftmargin=1.8em,itemsep=2pt,topsep=2pt]
\item 충분히 낮은 율의 준평형 코인 하프셀 $dQ/dV$ 를 여러 온도 $T_k$ 에서 측정한다(저율$\cdot$과평활 회피로
전이 폭이 동역학 lag 로 부풀지 않게 한다).
\item 측정 전위에서 분극을 먼저 분리한다 — Chapter 1 의 $V_n=V_\mathrm{app}-\sigma_d|I|R_n$ 로 IR$\cdot$과전압을
빼야 \emph{평형} $U_\oc$ 의 온도 의존만 남는다.
\item 같은 $Q_j$ 구조를 공유하되 각 온도의 중심 $U_j(T_k)$ 와 폭 $w_j(T_k)$ 를 허용해 Chapter 1 의 logistic
전이 모델~\eqref{eq:logistic} 을 \emph{동시} 피팅한다. 이는 반응별 $U_j^0(T)$ 를 온도 연속함수로 두는 MSMR
(Multi-Species, Multi-Reaction) 절차와 직접 대응한다 \cite{msmr_partI,msmr_partII}.
\item 회귀한 $U_j(T)$ 의 중심 기울기에서 $\Delta S^0_j=F\,\dd U_j/\dd T$ 를 얻고, Allart 등의 전극 분리
$\Delta S(x)$ \cite{allart2018} 와 부호$\cdot$규모를 대조해 검증한다(표~\ref{tab:ds}).
\item 봉우리 \emph{내부}의 $x$-의존은 따로 입력하지 않는다 — 식~\eqref{eq:single_config} 의 분포 config 항
$(R/F)\ln[\xi_j/(1-\xi_j)]$ 가 자동으로 채운다. 마지막으로 종합식과 $\dot Q_\rev=-IT\,\partial U_\oc/\partial T$
(Bernardi 출구 \cite{bernardi1985})로 가역 발열을 산출한다.
\end{enumerate}
\end{procedurebox}

% ====================================================================
\section*{맺음}
\addcontentsline{toc}{section}{맺음}
% ====================================================================
본 장은 코인 하프셀의 엔트로피와 가역 발열을 \emph{통계열역학}으로 세웠다. 분배함수 $Z$ 에서 점유 분포
$\avg{n}$ 을 끌어내 그것이 Chapter 1 logistic 의 기원임을 보였고($w=RT/F$ 가 분포 폭), 점유 분포에서
configurational 엔트로피와 그 발산하는 부분몰 항을, 같은 분포 언어로 vibrational(포논 BE)$\cdot$electronic
(FD Sommerfeld) 엔트로피를 세웠다. 세 분포가 합쳐 측정 $\partial U_\oc/\partial T(x)$ 의 모든 모양을 만든다:
겹침 가중(파생 A, 수치 검증 완료)으로 봉우리 사이 연속 블렌드를, 봉우리 내부 config(파생 B, 이중계산 분리)로
$x$-의존을, 상호작용 $w_\eff(\Omega)$(파생 C)로 폭 변형과 상분리 임계를, 히스 분기 평균(파생 D)으로 가역/비가역
분리를 닫았다. 극한 여섯 코너가 분포 식을 검산하고, ``전이당 상수 $\Delta S^0_j$ $+$ 분포''로 측정급 비선형이
자동 생성됨을 확정했다. 가역 발열 $\dot Q_\rev=-IT\,\partial U_\oc/\partial T$ 은 이 분포들을 재배열하는 열이다.

% ====================================================================
\begin{thebibliography}{99}
\bibitem{bernardi1985} D. Bernardi, E. Pawlikowski, J. Newman, ``A General Energy Balance for Battery Systems,'' \emph{J. Electrochem. Soc.} \textbf{132}(1), 5--12 (1985). DOI: 10.1149/1.2113792.
\bibitem{newman} J. Newman, K. E. Thomas-Alyea, \emph{Electrochemical Systems}, 3rd ed., Wiley (2004). (OCV($T$): slope $=\Delta S$, intercept $=\Delta H$; lattice-gas insertion thermodynamics.)
\bibitem{huggins2009} R. A. Huggins, \emph{Advanced Batteries: Materials Science Aspects}, Springer (2009). (격자기체 삽입 전극 열역학·Bragg--Williams 평형 전위.)
\bibitem{bazant2013} M. Z. Bazant, ``Theory of Chemical Kinetics and Charge Transfer Based on Nonequilibrium Thermodynamics,'' \emph{Acc. Chem. Res.} \textbf{46}(5), 1144--1160 (2013). DOI: 10.1021/ar300145c. (regular-solution lattice-gas OCV·상분리 임계 $\Omega=2RT$.)
\bibitem{reynier2003} Y. Reynier, R. Yazami, B. Fultz, ``The entropy and enthalpy of lithium intercalation into graphite,'' \emph{J. Power Sources} \textbf{119--121}, 850--855 (2003). DOI: 10.1016/S0378-7753(03)00285-4.
\bibitem{allart2018} D. Allart \emph{et al.}, ``Model of Lithium Intercalation into Graphite by Potentiometric Analysis with Equilibrium and Entropy Change Curves of the Graphite Electrode,'' \emph{J. Electrochem. Soc.} \textbf{165}(2), A380 (2018). DOI: 10.1149/2.1251802jes.
\bibitem{occupation2019} ``Transitions of lithium occupation in graphite: A physically informed model in the dilute lithium occupation limit supported by electrochemical and thermodynamic measurements,'' \emph{Electrochimica Acta} (2019). DOI: 10.1016/j.electacta.2019.135634. [방법 수준 인용 — 원문 식 미확보.]
\bibitem{chemmater2015} ``Intercalation Compounds from LiH and Graphite: Relative Stability of Metastable Stages\ldots,'' \emph{Chem. Mater.} \textbf{27} (2015). DOI: 10.1021/acs.chemmater.5b00235. [형성 ΔH calorimetry — abstract tier.]
\bibitem{jpcc2021} ``Thermodynamic Analysis of Li-Intercalated Graphite by First-Principles Calculations with Vibrational and Configurational Contributions,'' \emph{J. Phys. Chem. C} \textbf{125}(51) (2021). DOI: 10.1021/acs.jpcc.1c08992. [abstract tier.]
\bibitem{msmr_partI} ``Quantifying the Entropy and Enthalpy of Insertion Materials for Battery Applications Via the Multi-Species, Multi-Reaction Model,'' \emph{J. Electrochem. Soc.} (2024). DOI: 10.1149/1945-7111/ad1d27.
\bibitem{msmr_partII} ``Quantifying the Temperature Dependence of the MSMR Model Part II: Estimation of Entropy Coefficient for Meso-Carbon Micro-Bead Graphite,'' \emph{J. Electrochem. Soc.} (2024). DOI: 10.1149/1945-7111/ad70d9.
\bibitem{standardised2024} ``A Standardised Potentiometric Method for the Effective Parameterisation of Reversible Heating in a Lithium-Ion Cell,'' \emph{J. Electrochem. Soc.} (2024). DOI: 10.1149/1945-7111/ad4918.
\bibitem{hysteresis2018} ``Uncertainties in entropy due to temperature path dependent voltage hysteresis in Li-ion cells,'' \emph{J. Power Sources} (2018). DOI: 10.1016/j.jpowsour.2018.05.060.
\bibitem{numverif2026} 본 연구 내부 수치 검증(Derivative A), \texttt{Anode\_Fit\_v11} 4-전이 흑연 staging 파라미터, 유한차분 vs.\ 가중식 비교 (2026). [내부 자료 — 부동소수점 일치 PASS, 본문 \S\ref{ssec:overlap} 요지.]
\end{thebibliography}

\end{document}
